\section{The same question in time: the filter}
\label{sec:inffilter}

Here is the morning the gate cannot help with.  A one-month smile
has nine near-the-money quotes --- dense, liquid, well past every
requirement, so every gate of \cref{sec:infgate} is closed, and
rightly so.  But one of the nine is stale: its market moved and the
quote did not.  Chase the mids literally and the fit grows a
curvature kink that nobody traded.  This is not a gap --- the region
is superbly observed by count --- and using the prior to pull the
curvature back would violate the very guarantee
\cref{prop:infnodamp} establishes.  The problem is not absence of
data; it is \emph{noise} in data that is present.  The right
instrument treats today's reading as one more noisy estimate of a
moving quantity, to be weighed --- not vetoed --- against what
yesterday leads us to expect.  That instrument is the filter
\citep{Kalman1960,AndersonMoore1979}, and after
\cref{sec:infcombine} it can be assembled in a page.

\subsection{The state, and the cycle}

What should be filtered?  Not raw option mids: contracts appear and
disappear, the out-of-the-money side flips when the forward moves,
and a smoothed vector of raw volatilities has no arbitrage
guarantee.  The reference implementation filters the smile's three
\emph{ATM handles} --- the level $\sigma(0)$, the skew $s_0$
(Chapter~9's ATM slope), and the curvature (the second derivative of
the fitted volatility at the money) --- read off a fitted,
arbitrage-checked smile.  Three numbers, model-agnostic, stable in
meaning from day to day: a small state that summarizes what the
belly of the smile is doing.

Each handle then lives the same three-step cycle at every snapshot.
Write $O$ for the handle's true (unknown) value today.

\emph{Predict.}  Yesterday's combined estimate is pushed to today's
market.  Chapter~9 built exactly this transport: for a log-forward
move $H$ under the regime dial $\mathcal{R}$, the level moves by
$\mathcal{R}\,s_0H$, the skew by curvature times $H$, the curvature
is carried unchanged.  The transported value is the
\emph{prediction} $\hat O^{-}$, with a variance
$\mathcal{V}^{-}$ that must account for everything that can have
happened since yesterday (\cref{sec:infbudgets} builds it).

\emph{Observe.}  Today's prepared quotes are fitted --- data only,
no temporal prior --- and the handle is read off the fit: the
\emph{observation} $\hat O_{\mathrm{obs}}$, with a variance
$\mathcal{V}_{\mathrm{obs}}$ that must reflect how precisely, and
how consistently, today's quotes actually pin this handle
(\cref{sec:infbudgets} again).

\emph{Update.}  Two independent, unbiased estimates of the same
number, with variances: this is \cref{sec:infcombine} verbatim.
The precision-weighted average, rearranged into its standard form,
is the update:
\begin{equation}
\boxed{\;
\begin{aligned}
\mathcal{K}
&=\frac{\mathcal{V}^{-}}{\mathcal{V}^{-}+\mathcal{V}_{\mathrm{obs}}}
\;\in(0,1),\\
\hat O^{+}
&=\hat O^{-}+\mathcal{K}\,\bigl(\hat O_{\mathrm{obs}}-\hat O^{-}\bigr),\\
\mathcal{V}^{+}
&=(1-\mathcal{K})\,\mathcal{V}^{-}.
\end{aligned}
\;}
\label{eq:infupdate}
\end{equation}
The first line defines the \emph{gain} $\mathcal{K}$: the fraction
of the surprise the update accepts.  The second applies it to the
\emph{innovation} $\hat O_{\mathrm{obs}}-\hat O^{-}$ --- the gap
between what arrived and what was expected.  The third is the
posterior variance; it deserves its two lines of algebra.
Precisions add, \cref{eq:infprecadd}, so
\[
\frac{1}{\mathcal{V}^{+}}
=\frac{1}{\mathcal{V}^{-}}+\frac{1}{\mathcal{V}_{\mathrm{obs}}}
=\frac{\mathcal{V}^{-}+\mathcal{V}_{\mathrm{obs}}}
{\mathcal{V}^{-}\mathcal{V}_{\mathrm{obs}}}
\quad\Longrightarrow\quad
\mathcal{V}^{+}
=\frac{\mathcal{V}_{\mathrm{obs}}}
{\mathcal{V}^{-}+\mathcal{V}_{\mathrm{obs}}}\,\mathcal{V}^{-}
=(1-\mathcal{K})\,\mathcal{V}^{-},
\]
and the posterior mean, checked the same way, is
$\mathcal{V}^{+}\bigl(\hat O^{-}/\mathcal{V}^{-}
+\hat O_{\mathrm{obs}}/\mathcal{V}_{\mathrm{obs}}\bigr)
=\hat O^{-}+\mathcal{K}(\hat O_{\mathrm{obs}}-\hat O^{-})$.
Tomorrow, $\hat O^{+}$ and $\mathcal{V}^{+}$ become the input to the
next predict step, and the cycle repeats.

\subsection{The gain is computed, not set}

Pause on what \cref{eq:infupdate} does \emph{not} contain.  Every
fitter that carries state through time answers, somewhere, one
question: how far should this morning's quotes move yesterday's
surface?  The tempting implementation is a hand-set weight ---
tuned per desk, per mood.  In \cref{eq:infupdate} that weight is
the gain, and the gain is not a setting: it is an \emph{output},
the ratio of two variances.  The design question is therefore not
where to set the gain but how to build $\mathcal{V}^{-}$ and
$\mathcal{V}_{\mathrm{obs}}$ --- two quantities with definitions,
which \cref{sec:infbudgets} gives, and a test, which
\cref{sec:infaudit} gives.  When both variances match what they
measure, the gain follows the market: a noisy morning arrives with
a large $\mathcal{V}_{\mathrm{obs}}$ and is barely admitted; a
clean morning after a quiet night arrives against a small
$\mathcal{V}^{-}$ and is taken nearly whole.  The gate of
\cref{sec:infgate} and the gain of \cref{eq:infupdate} divide the
chapter's work: the gate answers ``\emph{where} did the market not
speak?''\ with a dead zone; the gain answers ``the market spoke
here --- \emph{how noisy was its voice?}''\ with a ratio that is
never zero and never one.

\subsection{The stale strike, filtered}

Now run the morning that opened this section through
\cref{eq:infupdate}, handle by handle.  The prediction ---
yesterday's state, transported --- says level $20.0\%$, skew
$-0.35$, curvature $0.10$, with standard deviations
$(0.30\%,\ 0.08,\ 0.05)$.  Today's data-only fit reads level
$20.4\%$, skew $-0.37$, curvature $0.55$: a plausible level move, a
plausible skew --- and a curvature five times the prediction,
manufactured by the one stale strike.  The observation variances,
built in the next section, come out as standard
deviations $(0.15\%,\ 0.05,\ 0.30)$: level and skew tight, curvature
wide, because the cluster that produced it contradicts itself.  Take
these six variances as given for one page; \emph{building} them is
\cref{sec:infbudgets}'s whole subject.

\begin{figure}[htbp]
\centering
\paperfig{\textwidth}{fig_flt_update.pdf}
\caption{The stale-strike morning through the scalar update
\cref{eq:infupdate}, one handle per row.  (a)~Each handle drawn on
its own normalized segment from prediction (circle, at 0) to
observation (square, at 1); the posterior (diamond) sits at exactly
its gain.  The level moves 80\% of the way to the market; the
curvature kink is admitted at less than 3\%.  (b)~The three computed
gains: ratios of stated variances, computed, not set.}
\label{fig:infupdate}
\end{figure}

\Cref{fig:infupdate}(a) draws each handle on a normalized segment:
prediction at zero, observation at one, so that --- by the second
line of \cref{eq:infupdate} --- the posterior's position \emph{is}
its gain, and \cref{eq:infupdate} can be checked by eye.  The level,
predicted to $0.30\%$ and observed to $0.15\%$, carries gain
$\mathcal{K}=0.09/(0.09+0.0225)=\MacFiltCaseGainLevel$ (the worked
number of \cref{sec:infcombine}): the posterior,
\MacFiltCasePostLevelPct\%, sits four-fifths of the way to the
market, with posterior standard deviation
\MacFiltCasePostSdLevelBp\ bp.  The skew behaves likewise at gain
\MacFiltCaseGainSkew\ (posterior \MacFiltCasePostSkew).  The
curvature is where the stale strike lands: prediction ten times tighter than the
observation, gain
$\mathcal{K}=0.05^2/(0.05^2+0.30^2)=\MacFiltCaseGainCurv$, posterior
\MacFiltCasePostCurv\ --- the five-fold kink is not vetoed by any
rule; it is \emph{priced}, and at that price the update accepts
under three percent of it.  Panel~(b) collects the three gains as
bars: computed per handle, unequal because the evidence is
unequal.

One remark on dimensions before moving on.  The three handles were
updated \emph{separately} --- three scalar filters side by side.  A
full three-dimensional update also exists, in which the handles'
estimated correlations let one handle's innovation move another's
posterior; it is the optimal version when the stated correlations
are true.  The danger is that on a coarse chain they are not:
with few strikes, level and curvature are estimated \emph{jointly}
rather than separately, and the correlation entries of the
observation covariance are then precisely its least reliable
numbers.  Spent at face value, they let a junk curvature innovation
drag the level along with it.  A per-handle update makes that
failure structurally impossible --- each scalar posterior lies
between its own prediction and its own observation, full stop.  The
reference implementation therefore combines per handle; the
discarded cross-handle information is real, and it is the cost of
not spending unreliable correlations.

Everything now rests on the two variances that fed
\cref{eq:infupdate}.  Where does
$\mathcal{V}_{\mathrm{obs}}$ come from, such that a
self-contradictory cluster shows up as a wide one?  And what makes
$\mathcal{V}^{-}$ grow at the right rate between snapshots?
